\documentclass[11pt,a4paper]{article}

% ── Codificación y fuentes ────────────────────────────────────────────────────
\usepackage[utf8]{inputenc}
\usepackage[T1]{fontenc}
\usepackage{lmodern}
\usepackage[spanish]{babel}

% ── Geometría y espaciado ─────────────────────────────────────────────────────
\usepackage[top=2.5cm, bottom=2.5cm, left=2.8cm, right=2.8cm]{geometry}
\usepackage{setspace}
\setstretch{1.15}
\usepackage{parskip}

% ── Tablas ────────────────────────────────────────────────────────────────────
\usepackage{booktabs}
\usepackage{tabularx}
\usepackage{array}
\usepackage{longtable}
\usepackage{enumitem}

% ── Colores y cajas ───────────────────────────────────────────────────────────
\usepackage[dvipsnames]{xcolor}
\usepackage{tcolorbox}
\tcbuselibrary{skins, breakable}

% ── Cabeceras y pies de página ────────────────────────────────────────────────
\usepackage{fancyhdr}
\usepackage{lastpage}
\pagestyle{fancy}
\fancyhf{}
\fancyhead[L]{\small\textbf{Informe de Evaluación} --- Física para Bioanalistas}
\fancyhead[R]{\small anamerc acuna 33279962}
\fancyfoot[C]{\small Página \thepage\ de \pageref{LastPage}}
\renewcommand{\headrulewidth}{0.4pt}

% ── Hiperenlaces ──────────────────────────────────────────────────────────────
\usepackage[colorlinks=true, linkcolor=NavyBlue, urlcolor=NavyBlue]{hyperref}

% ── Paleta de colores personalizados ─────────────────────────────────────────
\definecolor{desempeno_excelente}{HTML}{1A6B2F}
\definecolor{desempeno_bueno}{HTML}{2E5A9C}
\definecolor{desempeno_suficiente}{HTML}{8B6914}
\definecolor{desempeno_deficiente}{HTML}{C0392B}
\definecolor{desempeno_insuficiente}{HTML}{7B241C}
\definecolor{grisFondo}{HTML}{F5F5F5}
\definecolor{azulTitulo}{HTML}{1B2A6B}
\definecolor{verdeFortaleza}{HTML}{1A6B2F}
\definecolor{naranjaMejora}{HTML}{A04000}

% ── Estilos de cajas ─────────────────────────────────────────────────────────
\tcbset{
  cajaTitulo/.style={
    enhanced, breakable,
    colback=azulTitulo!8, colframe=azulTitulo,
    fonttitle=\bfseries\large, coltitle=white,
    attach boxed title to top left={yshift=-2mm, xshift=4mm},
    boxed title style={colback=azulTitulo},
    top=6pt, bottom=6pt, left=8pt, right=8pt,
  },
  cajaFortaleza/.style={
    enhanced, breakable,
    colback=verdeFortaleza!6, colframe=verdeFortaleza!60,
    fonttitle=\bfseries, coltitle=verdeFortaleza,
    top=4pt, bottom=4pt, left=8pt, right=8pt,
  },
  cajaMejora/.style={
    enhanced, breakable,
    colback=naranjaMejora!6, colframe=naranjaMejora!60,
    fonttitle=\bfseries, coltitle=naranjaMejora,
    top=4pt, bottom=4pt, left=8pt, right=8pt,
  },
  cajaRecomendacion/.style={
    enhanced, breakable,
    colback=grisFondo, colframe=gray!50,
    fonttitle=\bfseries, coltitle=black,
    top=4pt, bottom=4pt, left=8pt, right=8pt,
  },
}

% ─────────────────────────────────────────────────────────────────────────────
\begin{document}

% ══ PORTADA ══════════════════════════════════════════════════════════════════
\begin{center}
  {\Large\bfseries\color{azulTitulo} Universidad de Oriente/Núcleo Sucre --- Física para Bioanalistas}\\[4pt]
  {\large Informe de Evaluación Preliminar}\\[12pt]
  \rule{\linewidth}{1.2pt}\\[6pt]
  {\LARGE\bfseries anamerc acuna 33279962}\\[4pt]
  \rule{\linewidth}{0.6pt}
\end{center}

% ══ FICHA TÉCNICA ════════════════════════════════════════════════════════════
\vspace{4pt}
\begin{tcolorbox}[cajaTitulo, title=Datos del informe]
\begin{tabular}{@{}llll@{}}
  \textbf{Estudiante:}  & anamerc acuna 33279962  &
  \textbf{Fecha:}       & 2026-06-25 \\[4pt]
  \textbf{Archivo PDF:} & \texttt{anamerc\_acuna\_33279962.pdf}  &
  \textbf{Páginas:}     & 3 \\
\end{tabular}
\end{tcolorbox}

% ══ RESULTADO GLOBAL ═════════════════════════════════════════════════════════
\vspace{6pt}
\begin{center}
  \begin{tcolorbox}[
    enhanced,
    colback=white, colframe=desempeno_excelente,
    width=0.62\linewidth,
    boxrule=2pt,
    halign=center, valign=center,
    top=8pt, bottom=8pt,
  ]
    \centering
    {\large\bfseries Puntaje sugerido}\\[6pt]
    {\Huge\bfseries\color{desempeno_excelente} 9.0/10}\\[6pt]
    {\large Nivel de desempeño: \textbf{\color{desempeno_excelente} Excelente}}
  \end{tcolorbox}
\end{center}

% ══ RESUMEN GENERAL ══════════════════════════════════════════════════════════
\section*{Resumen general}
El estudiante demuestra un excelente dominio de los principios de la física de fluidos aplicados a problemas de bioanálisis. Ha resuelto correctamente la mayoría de los problemas, aplicando las fórmulas adecuadas y realizando los cálculos de manera precisa. Se observa una buena organización y claridad en la presentación de las respuestas, con un único error de magnitud en el cálculo final de uno de los problemas.

% ══ FORTALEZAS Y ÁREAS DE MEJORA ════════════════════════════════════════════
\vspace{4pt}
\begin{minipage}[t]{0.48\linewidth}
  \begin{tcolorbox}[cajaFortaleza, title={Fortalezas}]
\begin{itemize}[leftmargin=1.5em, itemsep=2pt]
  \item Excelente comprensión y aplicación de los principios de hidrostática, Arquímedes, continuidad y Bernoulli.
  \item Habilidad para identificar y utilizar las fórmulas correctas para cada problema.
  \item Correcta conversión de unidades y manejo de datos en la mayoría de los casos.
  \item Claridad y organización en el desarrollo de los problemas.
  \item Buen manejo de las leyes de Poiseuille y razonamiento proporcional.
\end{itemize}
  \end{tcolorbox}
\end{minipage}
\hfill
\begin{minipage}[t]{0.48\linewidth}
  \begin{tcolorbox}[cajaMejora, title={Areas de mejora}]
\begin{itemize}[leftmargin=1.5em, itemsep=2pt]
  \item Precisión en los cálculos numéricos finales, especialmente en problemas que involucran potencias de diez.
\end{itemize}
  \end{tcolorbox}
\end{minipage}

% ══ OBSERVACIONES POR PREGUNTA ═══════════════════════════════════════════════
\section*{Observaciones por pregunta}

\begin{longtable}{>{\bfseries}p{2.8cm} p{1.6cm} p{7.0cm} p{2.8cm}}
  \toprule
  \textbf{Pregunta} & \textbf{Puntaje} & \textbf{Evaluación} & \textbf{Errores detectados} \\
  \midrule
  \endhead
  \bottomrule
  \endfoot
    \textbf{Pregunta 1: Presión Hidrostática y Venosa} & 2.0/2.0 & La respuesta es completamente correcta. El estudiante aplicó correctamente el principio de la presión hidrostática para determinar la altura mínima requerida. Las unidades son consistentes y el cálculo es preciso. La segunda parte de la pregunta es ambigua y no se abordó, lo cual es razonable dada la falta de información adicional. & \textit{---} \\
    \midrule
    \textbf{Pregunta 2: Principio de Arquímedes} & 2.0/2.0 & La respuesta es completamente correcta. El estudiante calculó correctamente la fuerza de empuje, el volumen del tejido óseo y su densidad, aplicando el principio de Arquímedes y la definición de densidad. Los pasos son claros y los resultados precisos. & \textit{---} \\
    \midrule
    \textbf{Pregunta 3: Hemodinámica (Continuidad y Bernoulli)} & 2.0/2.0 & La respuesta es completamente correcta. El estudiante aplicó de manera impecable las ecuaciones de continuidad y Bernoulli para determinar la velocidad en la zona de estenosis y la presión resultante. Los cálculos son precisos y el razonamiento es sólido. & \textit{---} \\
    \midrule
    \textbf{Pregunta 4: Ley de Poiseuille (Análisis Proporcional)} & 2.0/2.0 & La respuesta es conceptual y procedimentalmente correcta. El estudiante aplicó la ley de Poiseuille para determinar la relación entre el caudal y el radio, calculando correctamente que el 40.96\% del caudal original se mantendrá. La interpretación de la pregunta ('que se mantendrá circulando') es plausible y se considera correcta en este contexto. & \textit{---} \\
    \midrule
    \textbf{Pregunta 5: Flujo Viscoso y Resistencia} & 1.0/2.0 & El estudiante identificó correctamente la ley de Poiseuille y la despejó para la diferencia de presión. La conversión de unidades (diámetro a radio) y la elevación a la cuarta potencia del radio son correctas. Sin embargo, el resultado numérico final presenta un error de magnitud significativo (un factor de 1000), lo que sugiere un error de cálculo o de manejo de potencias de diez en el paso final. & \textit{Error procedimental en el cálculo numérico final (magnitud del resultado)} \\
    \midrule
\end{longtable}

% ══ ERRORES DETECTADOS ═══════════════════════════════════════════════════════
\section*{Análisis de errores}

\subsection*{Errores conceptuales}
\textit{Ninguno identificado.}

\subsection*{Errores procedimentales}
\begin{itemize}[leftmargin=1.5em, itemsep=2pt]
  \item Error de cálculo numérico en la Pregunta 5, resultando en una diferencia de presión con una magnitud incorrecta (factor de 1000).
\end{itemize}

% ══ RECOMENDACIONES AL ESTUDIANTE ════════════════════════════════════════════
\section*{Retroalimentación para el estudiante}
\begin{tcolorbox}[cajaRecomendacion, title={Dirigido a: anamerc acuna 33279962}]
Estimado estudiante, tu desempeño en este examen es excelente. Demuestras una sólida comprensión de los principios fundamentales de la física de fluidos y su aplicación en el contexto del bioanálisis. Has resuelto la mayoría de los problemas de manera impecable, mostrando un buen manejo de fórmulas, unidades y razonamiento. Tu principal área de mejora es la precisión en los cálculos numéricos finales, especialmente cuando se trabaja con potencias de diez. Te sugiero revisar cuidadosamente los pasos de cálculo y el uso de la calculadora para evitar errores de magnitud. ¡Sigue así, tu base es muy fuerte!
\end{tcolorbox}

% ══ NOTA PARA EL PROFESOR ════════════════════════════════════════════════════
\section*{Nota para el profesor}
\begin{tcolorbox}[
  enhanced, breakable,
  colback=yellow!8, colframe=yellow!60!black,
  fonttitle=\bfseries, coltitle=black,
  title={Observaciones para revisión manual},
  top=4pt, bottom=4pt, left=8pt, right=8pt,
]
El estudiante ha demostrado un nivel de comprensión muy alto en todos los temas evaluados. Los problemas 1, 2, 3 y 4 fueron resueltos de manera impecable. En la Pregunta 4, la interpretación de 'porcentaje que se mantendrá circulando' fue tomada como el caudal final respecto al original, lo cual es una interpretación razonable de la ambigüedad de la pregunta. En la Pregunta 5, el planteamiento y las conversiones son correctas, pero el resultado final tiene un error de magnitud (factor de 1000). A pesar de este error puntual, el desempeño general es sobresaliente. El puntaje sugerido refleja la solidez conceptual y procedimental en la mayoría del examen.
\end{tcolorbox}

% ══ PIE DE INFORME ════════════════════════════════════════════════════════════
\vspace{12pt}
\begin{center}
  \footnotesize\color{gray}
  Informe generado automáticamente por el sistema de evaluación con IA (gemini-2.5-flash) y soporte LaTex. \\
  Este documento es una evaluación \textbf{preliminar} para ser revisado por el profesor antes de ser entregado al 
  estudiante. \\
  Generado el 2026-06-25 \textbullet\ BIOANALISIS Sistema Académico de Evaluación.
  \end{center}

\end{document}
